\section{学业与校园生活}\label{sec:academic}

丹麦高校强调自主学习、课堂参与和对学习过程负责。对来自中国的非欧盟留学生而言，学业问题还可能牵动学生居留许可：选课、考试、休学或转专业，不一定只是校内事务。最稳妥的做法是同时看三层规则：\textbf{国家法规与 SIRI 居留要求、所在院校的学习规则、具体课程与考试说明}。

\begin{grayleftbox}
\textbf{本章适用前提：}默认读者是以非欧盟身份持有丹麦学生居留许可的中国留学生。若你是交换生、联合培养学生、博士生，或持有其他类型居留许可，请向本校 International Office/Student Services 与 SIRI 单独核实。
\end{grayleftbox}

\subsection{先记住：学籍状态与学生居留许可相互关联}

SIRI 要求学生居留许可持有人积极参加获准就读的课程。是否属于“active student”通常由学校根据出勤、缺席及考试通过情况等作出判断；学校也有义务向 SIRI 报告不再积极就读的学生。若不再符合条件，SIRI 可能撤销或拒绝延长居留许可。

\begin{itemize}
  \item \textbf{休学（leave of absence）：}SIRI 明确说明，一般情况下休学期间不属于积极在读。不要只在学校系统提交休学申请；应先让学校国际学生办公室说明学籍影响，并向 SIRI 核实居留后果。
  \item \textbf{退学、被退学或提前结束项目：}学生居留许可所依据的学习目的可能不再成立。学校与 SIRI 的通知可能存在时间差，不要把“居留卡尚未过期”理解为仍可合法按原条件停留。
  \item \textbf{转专业、转校或开始新项目：}不要默认原许可自动覆盖新项目。在接受新录取、缴费或办理退学前，先让 SIRI 书面确认是否需要重新申请，以及何时申请。
  \item \textbf{延期毕业：}若学习进度延误并需要延长居留，应在原许可到期前办理；SIRI 要求申请人仍在同一项目注册并积极学习。学校通常需要提供在读与学习进度信息。
\end{itemize}

官方入口：\href{https://www.nyidanmark.dk/en-GB/Words-and-concepts/SIRI/Active-student}{SIRI：Active student}、\href{https://www.nyidanmark.dk/en-GB/You-want-to-extend/Study---extension/Higher-education}{SIRI：Extend a study permit}。居留问题应以 SIRI 的个案答复为准，而不是仅凭同学经验判断。

\subsection{看懂课程、学分与学习方式}

\subsubsection*{ECTS 与学习计划}

ECTS 是欧洲高等教育常用的学分单位。丹麦常见的完整本科为 180 ECTS、两年制硕士为 120 ECTS；也存在其他长度的项目，因此\textbf{本人项目的 curriculum/programme regulations 才是毕业要求的最终依据}。制定学习计划时，不要只把 ECTS 当作课时：它通常同时覆盖上课、阅读、作业、项目与备考。

每学期选课前，至少检查以下信息：

\begin{itemize}
  \item 课程是否属于必修、限选或自由选修，能否计入本人培养方案；
  \item prerequisite（先修要求）、academic prerequisites（考试资格要求）和 mandatory activities（强制活动）；
  \item 教学语言、课程时间、校区、容量限制及是否与其他课程或考试冲突；
  \item assessment form（考核形式）、允许携带的资料/设备、评分方式及重考安排；
  \item 选课、退课、考试报名和退考的截止日期。
\end{itemize}

\begin{grayleftbox}
\textbf{容易踩坑：}课程注册不一定等于考试注册；第一次考试自动注册，也不代表重考仍会自动注册。错过退考、未出席考试或未完成强制作业，在部分院校可能占用一次考试机会。请以本校当学年的 study rules 和学生系统状态为准，并保留成功注册或退选的确认邮件/截图。
\end{grayleftbox}

\subsubsection*{常见教学形式}

\noindent\begin{tabularx}{\textwidth}{@{}p{0.24\textwidth}X@{}}
\toprule
常见名称 & 你需要做什么 \\
\midrule
Lecture & 讲授课。课前阅读与课后练习常由学生自行安排，不一定逐项提醒。 \\
Exercise class/Tutorial & 习题课或辅导课，通常用于讨论方法与解题；部分课程会把参与或提交列为考试资格。 \\
Seminar & 研讨课，常要求阅读文献、发言、展示或互评。 \\
Lab/Workshop & 实验或工作坊，可能有安全培训、出勤、实验记录和报告要求。 \\
Project/Group work & 项目或小组作业。应尽早明确分工、版本管理、引用方式及个人贡献记录。 \\
Supervision & 与导师的指导会议。提前发送议程、材料和具体问题，会比临时泛泛讨论更有效。 \\
\bottomrule
\end{tabularx}

如果英文授课节奏较快，可在第一周就向教师确认阅读优先级、课程术语与作业标准。课堂讨论中的质疑通常针对观点而非个人；清楚说明“我目前的理解是……”并提出证据，是很常见的参与方式。

\subsection{考试、成绩与申诉}

\subsubsection*{先读考试说明，再按经验备考}

丹麦高校常见闭卷/开卷笔试、口试、居家作业、项目报告、展示及组合考核。同名课程在不同学期的考试形式也可能调整。请在备考前逐项确认：

\begin{itemize}
  \item 考试日期、地点、报到要求和迟到规则；
  \item 可用资料、计算器、电脑、网络与软件，以及是否允许生成式 AI；
  \item 是否有字数/页数限制、模板、匿名提交、文件格式与命名要求；
  \item 小组成果如何认定个人贡献，口试是否基于项目报告；
  \item 突发疾病、技术故障和缺考应联系谁、在多长时间内提交何种证明。
\end{itemize}

\subsubsection*{7 分制}

丹麦目前通行的 7-point grading scale 包括 12、10、7、4、02、00 和 -3。\textbf{02 是及格线，00 与 -3 不及格}；部分课程使用 Pass/Fail（通过/不通过）或 Approved/Not approved（批准/未批准）。成绩并非简单对应中国百分制，不建议自行做线性换算；申请其他学校或工作时，优先提交正式成绩单与学校提供的评分说明。

\begin{center}
\begin{tabular}{ccl}
\toprule
成绩 & 是否通过 & 简要含义 \\
\midrule
12 & 通过 & 表现优秀，只有少量不重要的不足 \\
10 & 通过 & 表现很好，但有若干小问题 \\
7  & 通过 & 表现良好，但存在一些不足 \\
4  & 通过 & 表现尚可，但有较明显不足 \\
02 & 通过 & 达到最低可接受程度 \\
00 & 不通过 & 未达到最低要求 \\
-3 & 不通过 & 完全不可接受 \\
\bottomrule
\end{tabular}
\end{center}

国家层面的高校法规入口可见丹麦高教与科研部的\href{https://ufm.dk/en/legislation/prevailing-laws-and-regulations/education/universities}{Universities: prevailing laws and regulations}，其中包括考试与评分相关法规。丹麦文原文具有法律效力，英文版本主要用于理解。

\subsubsection*{生病、特殊安排与申诉}

\begin{itemize}
  \item \textbf{考试生病：}不要带病硬撑或事后只向任课教师解释。立即按本校流程报病，并在规定期限内提交符合要求的证明；证明形式和费用由院校规则决定。
  \item \textbf{长期疾病、残障或学习困难：}尽早联系 student counselling、study administration 或 SPS/无障碍支持。国际学生可获得的具体支持及申请材料可能不同，不要等到考试周才提出。
  \item \textbf{考试申诉：}先索取评分依据或书面决定，区分对考试过程、法律问题和学术评价的异议。高校考试法规通常规定申诉期限为结果公布后的两周，但校内提交入口和计算方式应以本校通知为准。重评或重考的结果可能低于原成绩，提交前应阅读完整后果。
\end{itemize}

\subsection{学术诚信与生成式 AI}

“只是借鉴”“组员一起写的”或“AI 帮我润色”不自动代表合规。以下行为都可能触及学术不端：未标注地复制或改写他人内容、未经允许与他人合作、重复提交自己以前的作业、伪造或选择性处理数据，以及使用考试未授权的资料或工具。

\begin{grayleftbox}
\textbf{生成式 AI 的安全判断顺序：}
\begin{enumerate}
  \item 先看具体课程页面、考试说明和本专业/学院规则是否允许；“学校总体鼓励 AI”不等于本次考试允许。
  \item 若规则不清楚，在提交前用学校邮箱向课程负责人提出具体问题，并保存书面答复。
  \item 允许使用时，按要求披露工具、版本、用途和对内容的影响；保留提示词、输出、修改记录与原始资料。
  \item 不向公共 AI 工具上传个人敏感信息、未公开研究数据、受保密协议约束的材料或他人的作品。
  \item 最终责任仍由学生承担：逐条核实事实、引用、代码、计算与参考文献，确保自己能够解释并捍卫提交内容。
\end{enumerate}
\end{grayleftbox}

不同学校甚至不同课程的界线会不同。例如，\href{https://fbu.ku.dk/english/generative-ai/}{哥本哈根大学关于生成式 AI 与良好学术实践的说明}强调诚信、透明与学生对最终成果负责；\href{https://www.ai.dtu.dk/rules/}{DTU 的 AI 规则入口}则提醒学生检查课程和考试的具体许可。它们可以帮助理解原则，但不能替代你所在院校和本次考试的规则。

小组作业建议建立简单的贡献记录：会议纪要、分工表、共享文档版本历史、代码提交记录和所用数据来源。这样既方便协作，也能在教师询问时说明个人贡献。

\subsection{学业波动时，学校与居留两条线同时处理}

\noindent\begin{tabularx}{\textwidth}{@{}p{0.18\textwidth}X X@{}}
\toprule
情况 & 校内优先动作 & 非欧盟学生额外动作 \\
\midrule
挂科或进度落后 & 联系 study adviser，核对重考、考试次数、学习活动要求与新学习计划。 & 检查是否仍会被学校认定为积极在读，以及居留到期日是否覆盖新的完成时间。 \\
考试生病 & 按时办理 sick exam/缺考证明，确认是否占用考试机会。 & 若疾病将导致长期中断或延期，向学校国际办公室与 SIRI 核实影响。 \\
考虑休学 & 先了解批准条件、复学时间、学费及校内账户权限。 & 提交申请前核实居留是否会被撤销、何时需要离境，以及复学时是否要重新申请。 \\
转专业或转校 & 获取新录取、学分转换结果和原项目结束日期。 & 让 SIRI 确认旧许可是否覆盖过渡期、是否需要新许可；不要自行假定。 \\
准备退学 & 索取书面确认、成绩单及已修学分记录。 & 立即了解合法停留基础、离境期限、住址与其他登记后续。 \\
\bottomrule
\end{tabularx}

与学校或 SIRI 沟通时，建议一次写清姓名、出生日期/案号（如适用）、学校与项目、当前学籍状态、关键日期、计划采取的动作和需要确认的问题。涉及个人信息时只使用官方渠道，不在公开群聊发布护照、居留卡或完整案号。

\subsection{新生第一个月学业清单}

\begin{enumerate}
  \item 激活学校邮箱、学生门户、学习平台和多重验证；学校正式通知通常只发到校内邮箱。
  \item 下载本人 curriculum/programme regulations 和当学年 study rules，记录文件版本与日期。
  \item 在个人日历中录入选课、退课、考试报名、退考、作业和重考截止日期。
  \item 逐门保存课程页面、learning objectives、考试形式、允许工具与 AI 说明。
  \item 核对学生系统中的课程与考试注册状态；发现异常立即用书面方式联系 administration。
  \item 找到 study adviser、International Office、student counselling、IT support 和校内紧急联系入口。
  \item 建立文件备份与引用管理习惯；重要提交保留最终文件、回执与时间戳。
  \item 检查护照和居留许可有效期，在到期前预留学校证明与 SIRI 申请时间。
\end{enumerate}

\begin{grayleftbox}
\textbf{信息核验说明：}本章国家法规与 SIRI 相关信息核验至 2026 年 8 月 12 日。院校的选课窗口、考试次数、病假证明、AI 使用和申诉入口可能每学年调整，请以学校学生门户及具体课程当期说明为准。
\end{grayleftbox}
